\subsection{Nonzero denominators before any sign assumption}
The weakened norm no longer has a constant one on its right-hand side.
It therefore does not, by itself, imply $M\ne0$. The correct order is
as follows. The affine oddness row gives $q\ne0$. Apply
Lemma~\ref{lem:signed-power} immediately to the request for $M$ at base
$q$. For $|q|\ge2$ it gives $x\ge0$ and $M=q^x$; for $q=\pm1$
it gives $M=q^x$ at every signed index. Either way $M\ne0$.
The request at the already positive index $t$ gives
$V=(F^2+1)^t\ge1$, including $F=0$.

From $Bv=M$ we obtain $Bv\ne0$. Set $b=|B|$ and $s=BF-V$.
The other new row gives $S=vs$, and the sphere gives
\begin{equation}\label{eq:wide-near}
 |s|\le b/\sqrt2<3b/4.
\end{equation}
If $F=0$, then $V=1$ and $s=-1$, so $b\ge2$. If
$f=|F|\ge1$, opposite signs of $B,F$ would force
$V=BF-s<-b+3b/4<0$. Thus $BF=bf>0$. Moreover
$V<b(f+3/4)$. For every integer $f\ge1$,
\[
 (f^2+1)^t\ge(f^2+1)^2>(f+1)(f+3/4),
\]
since the last difference is $f^4+f^2-7f/4+1/4>0$.
It follows that $b\ge f+2\ge3$, and therefore
\begin{equation}\label{eq:early-wide}
 V<b(f+3/4)\le b(b-5/4)<b^2-b<b^2.
\end{equation}
In particular every branch has $b\ge2$.

If $q=\pm1$, its power request gives $|M|=1$, contradicting
$Bv=M$. Hence $Q:=|q|\ge3$, and $|q+1|\ge2$.
Every actual representative $a_\nu$ is an exponent requested at base
$q+1$, so Lemma~\ref{lem:signed-power} gives $a_\nu\ge0$ and
$H\ge1$. The same argument gives $h\ge0$, while the semantic anchor
has already excluded $h=0$. The request at base $q$ now gives
$x\ge0$ and $M=q^x$. The retained denominator-ratio equation
$q^{t+1}v_{t+1}=M$, interpreted as exact nonzero powers and taken in
absolute value, forces $x\ge t+1\ge17$. Consequently the self-radix
request has base $BH$ of magnitude at least two and gives
\begin{equation}\label{eq:wide-radix}
 q=(BH)^x,\qquad Q=(bH)^x\ge b^x\ge b^{t+1}.
\end{equation}
If $F=0$, we have so far proved only $b\ge2$. But $q$ is odd and
$x>0$, so $b$ is odd; hence $b\ge3$. This parity step is necessary
when widening the former radius $1/2$.

\subsection{Exact factorial digits at a signed radix}
For each normalized row, $L_i v_i=M\ne0$ and
$v_iU_i=MQ_i+R_i$. Hence $v_i\mid R_i$, and $r_i=R_i/v_i$ is
an integer satisfying
\begin{equation}\label{eq:wide-remainder}
 U_i=L_iQ_i+r_i,\qquad |r_i|\le |L_i|/\sqrt2<3|L_i|/4.
\end{equation}
This uses no positivity of $L_i$ or $r_i$.

For the factorial digit, set $n=F^2+t$ and
$C_j=\binom nj$. Its exact numerator request gives $U=(q+1)^n$,
while its restored digit row gives
\[
 U=q^t(B+q\eta)+r.
\]
The signed low polynomial
$\rho=\sum_{j=0}^{t-1}C_jq^j$ satisfies
\[
 |\rho|<Q^t\frac{\max_{j<t}C_j}{Q-1}.
\]
For $F\ne0$, the elementary coefficient estimate
$C_j\le(F^2+1)^t=V$ holds for all $j\le t$. To prove its full
range, induct on $t$ in
$\binom{m+t}{j}\le(m+1)^t$, for $m\ge1$ and $0\le j\le t$.
Pascal's identity handles the next row through $j=t$, since
$2\le m+1$. Its new edge is bounded by
\[
 \binom{m+t+1}{t+1}=\prod_{i=1}^{t+1}(1+m/i)\le(m+1)^{t+1}.
\]
The initial case $t=0$ is immediate. Equations
\eqref{eq:early-wide}--\eqref{eq:wide-radix} now give
\[
 |\rho|<Q^t\frac{V}{Q-1}
       <Q^t\frac{2}{b^{t-1}}<Q^t/32.
\]
For $F=0$ one must not use $V=1$ as a coefficient bound.
Instead $C_j=\binom tj\le2^t$ and $b\ge3$, so
\[
 |\rho|<Q^t\frac{2^t}{3^{t+1}-1}
       <Q^t(2/3)^{t+1}<Q^t/32.
\]
The difference $r-\rho$ is divisible by the possibly negative integer
$q^t$, but its absolute value is less than
$(3/4+1/32)Q^t=25Q^t/32<Q^t$. Thus $r=\rho$.
Subtracting, dividing by $q^t$ and reducing modulo $q$ gives
$B\equiv C_t\pmod q$. For $F\ne0$,
$|B-C_t|\le b+C_t<b+b^2<Q$; for $F=0$, $C_t=1$ and
$|B-1|\le b+1<Q$. In both cases,
\begin{equation}\label{eq:wide-binomial}
 B=\binom{F^2+t}{t}>0.
\end{equation}
The branch $F=0$ would force $B=1$, contrary to $b\ge3$.
Thus $F\ne0$, the earlier common-sign argument gives $F>0$, and
\eqref{eq:wide-radix} finally gives $q>0$.
The negative-radix branch is excluded by the extracted arithmetic,
not by an initial sign assumption. All other requests, including
$N=B^x$, now have their ordinary positive meanings.

\subsection{A factorial fixed point with radius three quarters}
\begin{lemma}\label{lem:wide-factorial}
For integers $t\ge16$ and $F\ge1$, let
\[
 R_t(F)=\frac{(F^2+1)^t}{\binom{F^2+t}{t}}.
\]
Then $|F-R_t(F)|<3/4$ if and only if $F=t!$. At that solution,
$0<F-R_t(F)<1/32$.
\end{lemma}
\begin{proof}
Write $K=t!$ and
\[
 R_t(F)=K\prod_{i=1}^t\frac{F^2+1}{F^2+i}.
\]
Every factor is nondecreasing with $F$, and
\[
 R_t(F)\ge R_t(1)=\frac{2^t}{t+1}>2t^2+3/4.
\]
The last inequality holds at $16$ and propagates by induction: if
$A_t=(t+1)(2t^2+3/4)$, then
$2A_t-A_{t+1}=2t^3-4t^2-37t/4-4>0$ for $t\ge16$.
The distance hypothesis therefore gives $F>2t^2$, whence
$F\ge2t^2+1$. Since $R_t(F)<K$, the same hypothesis implies the
integer bound $F\le K$.

Furthermore
$F^2/(F+1)=F-1+1/(F+1)>2t(t-1)$, so
$F^2>2t(t-1)(F+1)$. For $u_i=(i-1)/(F^2+i)$,
\[
 \sum_i u_i<\frac{t(t-1)}{2F^2}<\frac1{4(F+1)}.
\]
Using $\prod_i(1-u_i)\ge1-\sum_i u_i$ gives
$R_t(F)>K(1-1/[4(F+1)])$. If $K\ge F+1$, this is
strictly greater than $F+3/4$, contradicting the distance hypothesis.
Thus $F=K$.

Conversely, at $F=K$ the same product estimate gives
\[
 0<K-R_t(K)<\frac{t(t-1)}{2K}<\frac1{32}.
\]
The final inequality holds at $16$ and its left side decreases thereafter.
This proves both the equivalence and the stronger canonical margin.
\end{proof}

By \eqref{eq:wide-binomial}, estimate \eqref{eq:wide-near} is exactly
$|F-R_t(F)|<3/4$. Hence $F=t!$ and $0<s=BF-V<B/32$.
The old ceiling witnesses $s$ and $B-s-1$ are nonnegative integers.
Also $q\ge B^{t+1}>B$, so $q-B-1\ge0$. Fresh ternary $G$ witnesses
restore all three deleted initial scalar sign certificates.

\subsection{All remaining digits and the refutation}
Only now use the true positive factorial front in the beta endpoints.
Since $y=V(V-1)>F+2$, the first endpoint with $u_0\le-1$ would give
$x\le F-y<0$. Thus $u_0\ge0$ and $x\ge F+1\ge2$.
The last endpoint then forces $u_t\ge1$, and
\begin{equation}\label{eq:wide-endpoint}
 x\ge2+(t+1)y.
\end{equation}
All four denominator levels are now positive divisors of $M=q^x$.
The strict amplification margin is unchanged. With
$\mathcal R=2+t+x+y+F$, one has $\mathcal R<2x$ and
$x\ge16!+1>2^{32}$; therefore
\[
 \mathcal R^{2^{23}}<(2x)^{2^{23}}<2^x.
\]
For the last inequality, the function $x/\log_2(2x)$ increases for
$x\ge2^{32}$, and at $2^{32}$ it exceeds $2^{23}$.
This is precisely the retained local precision margin.

The encoded-index congruence is $B\mid C_0+F^2+1$.
Since $B>F^2+1$ and $C_0\ge0$, we have $C_0>0$ and $H\ge2$.
Each ordinary representative $a\ge0$ has
$\binom aj\le\max(1,a)^k\le H^x$ at its range position $k\le x$,
including $a=j=0$. The retained constructor degree ledger gives
$h\le H^{2^{22}}$, while the factorial and endpoint inequalities give
$2^{22}(t+1)<x$. Thus all needed derived-label coefficients also are
at most $H^x$. Since $q=B^xH^x$, every genuine remaining low tail obeys
\[
 0\le\rho_i<L_i\frac{H^x}{q-1}
              <L_i\frac2{B^x}<L_i/32.
\]
Equation~\eqref{eq:wide-remainder} and the same divisible-difference
argument force $r_i=\rho_i$. Every original interval slack
$L_i-r_i-1$ is nonnegative and admits its old guard witnesses.
The quotient congruences are exactly the retained $N$-divisibility range
tests, with $N=B^x\mid q$. The unchanged prime-power extraction,
constructor equations and Horn-flow theorem therefore recover a finite
C222 refutation, and the named calibration gives a finite ZFC refutation.
This recovers a complete calibrated witness, not just some necessary
arithmetic conditions. Stored representative-minus-one coordinates may
still be $-1$; only the actual semantic representatives must be nonnegative.

\subsection{Constructive completeness for the wider sphere}
Start with a finite calibrated refutation. The retained cofinal construction
chooses a sufficiently large $t=16+d^2$, sets $F=t!$,
$V=(F^2+1)^t$, $y=V(V-1)$, and chooses a sufficiently large beta-CRT
code $x$. The true beta quotients are nonnegative and less than $x$,
since their denominators are at least two. All original strict range
witnesses therefore exist.

Every factor $1+F^2/i$ is odd because $F$ is even and $i\mid F$.
Thus $B$, $N=B^x$ and the precision modulus $N^2$ are odd. Adding
$N^2$ to a representative preserves its CRT residues and reverses parity.
Choose the dominant representative $a$ in $[2N^2,4N^2)$ with parity
opposite to $b$, retaining the others in $(0,2N^2]$. The literal
constructors give $h\ge2a^2$ and $h\equiv1+a+b\pmod2$.
Consequently $h$ is even and remains a dominant semantic anchor.
Choose a third independent representative's parity so the sum of all
nine is even, making $H$ odd. These choices preserve all precision
bounds. Then $q=(BH)^x$ is odd, and $z=(q-1)/2$ is an integer.
Affine oddness requires no parity condition on $x$.

For the deterministic comparison center take an admissible positive $A$,
$C=Y_h(A)$ and $U=X_h(A)$. Since $h$ is even, $C$ is even and $U$ is
odd. Set $l=hC$. Pell composition gives $J_{\rm p}=Y_l(A)=CY_C(U)$.
The Pell binomial expansion and $U^2-1=(A^2-1)C^2$ imply
\[
 \frac{Y_C(U)}C\equiv U^{C-1}\pmod C.
\]
Indeed the zeroth term after division by $C$ is $U^{C-1}$; every
higher term contains a factor $(U^2-1)^j/C$, hence at least one
factor $C$. Thus $J_{\rm p}/C^2$ is an odd integer and supplies $e$.
Set $H_{\rm p}=X_l(A)$, $f=(1-A)H_{\rm p}$ and
$D=A+H_{\rm p}f$. The pair $(X_h(D),Y_h(D))$ supplies both required
comparison residues, including when $D<0$. The retained cofinal frame
and cap construction provides its other witnesses. Canonical pairs have
residues $(1,k)$ modulo $A-1$, so the full orientation pool holds.

Finally $B\mid q\mid M$. Set $v=M/B$ and $S=v(BF-V)$.
Lemma~\ref{lem:wide-factorial} gives $0<S<M/32$.
For each true remainder take $R_i=(M/L_i)r_i$, with
$0\le R_i<M/32$. Thus
\[
 2\left(S^2+\sum_{i=1}^{11}R_i^2\right)
       <\frac{24}{1024}M^2=\frac3{128}M^2<M^2.
\]
The nonnegative integer difference has a $G$ representation, supplying
the new sphere row. Every other source witness is retained.
This constructs a new zero from every finite refutation; it is not a
bijection between arbitrary old and new ternary-guard assignments.

\subsection{Exact algebra after the arithmetic change}
The remaining reductions use three finite operations. An integral affine
coordinate shear has its explicit integral inverse. An elementary row
change $R_i\leftarrow R_i+gR_j$, $i\ne j$, preserves simultaneous zeros
and is inverted with $-g$. A unit-pivot row
$\epsilon x_v+f(\widehat{x_v})=0$, $\epsilon=\pm1$, uniquely restores
$x_v=-\epsilon f$ and can be deleted after that substitution in every
other row. The certificate verifies that the proposed pivot is independent
of $x_v$, that every transformed row remains quadratic when claimed,
and that all coordinate indices are compacted consistently.

These transformations preserve zero existence but need not preserve the
value of the sum of squares away from zeros. The final polynomial is the
sum of squares of the final literal residuals, not an asserted
coefficientwise identity with an earlier quartic. Combining both arithmetic
witness directions with the exact finite transformations gives
\eqref{eq:calibration} for every selected endpoint; negation gives
\eqref{eq:consistency}. The full source calibration is retained, not
replaced by the finite coefficient checker.
